\chapter{Background and Related Work}
\label{ch:background}

\section{Spiking neural networks}

Spiking neural networks (SNNs) are often described as the third generation of
neural networks~\cite{maass1997networks}. They take their inspiration from biological neurons, which communicate through
brief, all-or-nothing electrical impulses called \emph{spikes}.
In a conventional
ANN, each neuron outputs a continuous value computed once per input, while in an SNN,
information is carried by sequences of binary events distributed over time: at
each of $T$ discrete time steps, a neuron either emits a spike or stays silent.
The network therefore processes an input as a short temporal sequence. This holds
even when the input is static, as with a spectrogram, our specific case, which is presented
repeatedly or encoded as spikes over the $T$ steps.

This representation is exactly what makes SNNs attractive for low-power computing. Because
activations are binary, the multiply--accumulate operations that dominate the cost
of an ANN reduce to simple accumulations, performed only when a spike actually
is fired. Computation is therefore sparse and event-driven, with a silent neuron costing
almost nothing. Neuromorphic processors such as Intel's
Loihi~\cite{davies2018loihi} are built to exploit this sparsity directly, and the
energy advantage of an SNN is set to grow as its firing activity decreases.

The most widely used neuron model in deep SNNs is the leaky integrate-and-fire
(LIF) neuron~\cite{gerstner2014neuronal}. In this configuration, neuron maintains an internal state,
its membrane potential $U$, which
integrates incoming input, decays (``leaks'') towards rest over time, and triggers
a spike when it crosses a threshold $\theta$.

As shown in Figure~\ref{fig:neuron_models}, the LIF model sits between
biologically detailed neuron models and the artificial neurons of
conventional deep learning (ANNs).

\begin{figure}[htbp]
  \centering
  \includegraphics[width=\linewidth]{lif_neuron_models.png}
  \caption[Neuron models: from Hodgkin--Huxley to the artificial neuron]{%
    Neuron models ordered from biological realism to computational
    utility. Left: the Hodgkin--Huxley model describes the generation of
    action potentials through voltage-dependent ionic conductances, at a
    high computational cost. Centre: the leaky integrate-and-fire (LIF)
    neuron reduces the membrane to an RC circuit: the capacitor $C$
    integrates the input current $I_{\mathrm{in}}$, while the resistor $R$ makes
    the stored charge leak away over time, and a spike is emitted when the
    membrane potential reaches the threshold (denoted $\vartheta$ in the
    figure and $\theta$ in the text), after which the potential is reset.
    Unlike the Hodgkin--Huxley model, the membrane potential of a LIF neuron is an internal
    state and the spike is the output. Right: the artificial neuron of
    conventional ANNs, which computes a weighted sum of its inputs followed by an
    activation function, with no temporal dynamics and no spikes.
    Reproduced from the snnTorch documentation~\cite{snntorch_tutorial2}
    (MIT License).}
  \label{fig:neuron_models}
\end{figure}

In the discrete-time formulation used by snnTorch~\cite{eshraghian2023}, the
dynamics of a neuron at time step $t$ are the following
\begin{equation}
  U[t] = \beta\, U[t-1] + W X[t] - S[t-1]\,\theta ,
  \label{eq:lif}
\end{equation}
\begin{equation}
  S[t] =
  \begin{cases}
    1 & \text{if } U[t] > \theta, \\
    0 & \text{otherwise},
  \end{cases}
  \label{eq:spike}
\end{equation}
where $W X[t]$ is the weighted input from the previous layer, $S[t]$ is the binary
output spike, and the subtraction term resets the potential by $\theta$ after each
spike. The decay factor $\beta \in [0,1)$ sets how much of its past state the
neuron retains: values close to $1$ give a long memory, so the neuron integrates
evidence over many time steps; smaller values make it respond mainly to recent
input.

The property of the membrane potential to give the LIF neuron a memory is what separates
it from a stateless activation function: the output at time $t$ depends on the
entire history of inputs up to $t$.
The same mechanism, nonetheless, also creates the main
difficulty in training. The spike function in Eq.~\eqref{eq:spike} is a step
function, whose derivative is zero almost everywhere and undefined at the
threshold, so standard backpropagation cannot pass through it. The usual solution
is the \emph{surrogate gradient} approach~\cite{neftci2019surrogate}: the step
function is kept in the forward pass, but in the backward pass its
derivative is replaced with that of a smooth approximation, such as an arctangent.
Combined with backpropagation through time over the $T$ steps, this allows deep SNNs to
be trained with the same gradient-based tools as conventional networks.

These properties change the quantization problem in three main direct ways. First, SNN
activations are already binary, so the precision that quantization removes lies
in the weights and in the membrane potential, a state variable with no direct
equivalent in an ANN. Second, the threshold makes the network highly sensitive to
small numerical changes. Indeed, a perturbation that moves $U$ slightly above or below
$\theta$ turns a spike into silence, or the reverse. Because the membrane carries
its state forward, such an error does not stay local, but can propagate through
later time steps and layers. Third, explanations of an SNN can be temporal as well
as spatial. The Spike Activation Map (SAM)~\cite{kim2021sam} builds
attribution maps directly from the timing of spikes, without computing
gradients, and therefore show not only \emph{where} the network finds evidence in
the input but also \emph{when} during the $T$ time steps it does so. Whether
quantization preserves this temporal reasoning, and not only accuracy, is the
question this thesis sets out to answer.

\subsection{Training by surrogate gradients}
\label{sec:surrogate}

Unrolled over its $T$ time steps, a network of LIF neurons is effectively a recurrent
network: the membrane potential is a hidden state that each step reads and
updates (Equation~\ref{eq:lif}). Training it by gradient descent requires solving
two credit-assignment problems at once~\cite{neftci2019surrogate}: a
\emph{spatial} one, across layers, as in any multi-layer network, and a
\emph{temporal} one, since a weight influences the loss through every later
state of the neurons it feeds. Backpropagation through time (BPTT) solves both
by applying the chain rule to the unrolled graph. For a weight $W$ feeding a
neuron with potential $U[t]$ and output $S[t]$,
\begin{equation}
\frac{\partial \mathcal{L}}{\partial W}
= \sum_{t} \frac{\partial \mathcal{L}}{\partial S[t]}\,
  \frac{\partial S[t]}{\partial U[t]}\,
  \frac{\partial U[t]}{\partial W},
\qquad
\frac{\partial S[t]}{\partial U[t]} = \Theta'\!\left(U[t] - \theta\right),
\label{eq:bptt}
\end{equation}
where $\Theta$ is the Heaviside step of Equation~\ref{eq:spike} and
$\partial U[t]/\partial W$ is the total derivative through all earlier states.
The price is
memory: every state of every time step must be kept until the backward pass,
a cost that grows as $O(NT)$ for a layer of $N$ neurons.

\paragraph{The obstacle.} Every path through Equation~\ref{eq:bptt} contains the
factor $\Theta'$, which is zero everywhere except at the threshold, where it is
undefined. The gradient is therefore either zero or meaningless: a small change
of a weight almost never flips a spike, so the loss is piecewise constant in the
weights, with flat plateaus on which gradient descent cannot move. This is the
same problem met by networks with binary activations, and the solutions
developed for spiking networks are inspired by those~\cite{neftci2019surrogate}.

\paragraph{The surrogate.} A surrogate gradient leaves the model unchanged and
alters only its derivative. The forward pass keeps the exact binary spike, while in
the backward pass every occurrence of $\Theta'$ is replaced by the derivative
$\sigma'$ of a smooth function peaked at the threshold,
\begin{equation}
\frac{\partial S[t]}{\partial U[t]} \;\longrightarrow\;
\sigma'\!\left(U[t] - \theta\right).
\end{equation}
Neurons close to threshold, meaning those for which a small change could flip a spike,
receive the largest gradient, and gradient flows through silent neurons as well.
The resulting update is the gradient of a
\emph{virtual} surrogate loss, never computed explicitly, which follows the true
one while remaining continuous and finite~\cite{neftci2019surrogate}. This
distinguishes the approach from \emph{smoothed} spiking networks, which involve replacing
the spike itself with a soft or probabilistic function and so change the model
that is deployed.

\paragraph{The choice of surrogate.} Piecewise-linear, fast-sigmoid and
exponential surrogates have all been used successfully, all three sharing non-linearity 
and monotonically increasing towards the threshold. The
diversity of working choices suggests that the method does not depend
critically on the exact shape~\cite{neftci2019surrogate}. This work uses the
arctangent surrogate~\cite{fang2021plif} as implemented in
snnTorch~\cite{eshraghian2023},
\begin{equation}
\sigma'(x) = \frac{\alpha/2}{1 + \left(\frac{\pi}{2}\,\alpha\, x\right)^{2}},
\qquad \alpha = 2,
\end{equation}
which equals $1$ at the threshold and falls off as the inverse square of the
distance from it, so that neurons far from threshold still receive a small,
non-zero gradient. The reset term $-S[t-1]\,\theta$ of
Equation~\ref{eq:lif} is treated as a constant in the backward pass: the
recurrence then contributes $\partial U[t]/\partial U[t-1] = \beta$ alone, without
a second surrogate factor through the reset, a simplification that has been
found empirically to improve training~\cite{neftci2019surrogate}. With this training
configuration, spiking networks can manage to reach accuracies comparable to conventional architectures
on image benchmarks~\cite{eshraghian2023}.

The same device reappears in Section~\ref{sec:quantization-background}: the
straight-through estimator used to train through a rounding
operation~\cite{bengio2013ste} replaces the zero derivative of a step function
with that of a smooth one, exactly as the surrogate gradient does for the spike.

\paragraph{A consequence for explainability.} The substitution motivates the
choice of explanation method in Section~\ref{sec:sam-background}: any
gradient-based attribution computed on such a network inherits the
approximation error of the surrogate, because the gradients it uses are not the
gradients of the function the network actually computes.

\subsection{Input encoding}
\label{sec:encoding}

A spiking network processes its input over $T$ discrete time steps, while an
image is a single static frame, so the frame must be turned into a sequence. In this context, two
strategies dominate. In \emph{rate coding} each normalized pixel intensity is
used as the probability of emitting a spike at every step, and the value is
carried by the spike count; in \emph{direct} (or \emph{constant}) encoding, the
analog frame itself is presented, identical, at every time step.

Rate coding trades its simplicity in time steps: a pixel is represented by
counting spikes, so with $T$ steps it can take only $T + 1$ distinct values, and
the sampling makes the representation stochastic: the discretization error
shrinks only as $T$ grows, which is why rate-coded and converted networks
typically require between $100$ and $2\,000$ time steps to approach the
accuracy of their non-spiking counterparts~\cite{rathi2023dietsnn}.

Direct encoding, as introduced by DIET-SNN~\cite{rathi2023dietsnn}, removes
that error at the source. ``The pixel intensities of an image are applied
directly to the input layer of the SNN at each timestep'', and the first
convolutional layer, followed by LIF neurons, becomes at once a feature
extractor and a learned spike generator: it integrates the weighted pixel
values and fires when its membrane potential crosses threshold. Since the
input carries its full precision at every step, the number of steps is no longer
bounding the precision of the input. DIET-SNN combines this encoding with a
trained membrane leak and firing threshold per layer, which increase activation
sparsity, and reaches $69\,\%$ top-1 accuracy on ImageNet with $5$ time steps
and about $12\times$ less compute energy than the equivalent conventional
network~\cite{rathi2023dietsnn}. This work adopts the encoding but not the
trained neuron parameters: $\beta$ and the threshold are shared by all layers,
and $\beta$ is chosen by the constrained search of
Section~\ref{sec:constrained-search}.

Direct encoding has three remarkable consequences that the rest of this thesis relies on.

\paragraph{Temporal structure is generated inside the network.} Since the
stimulus does not vary, \textbf{all temporal structure in the network
originates from the internal LIF dynamics} --- leak, integration and reset ---
and not from variation in the input. The time axis of a spike train, and
therefore of a Spike Activation Map, is the time the network takes to integrate
evidence, which is not the time of the glitch: in an Omega-scan spectrogram, physical
time runs along the horizontal axis of the image and is, for the network, a
spatial dimension. Every temporal analysis in this thesis is to be read in that
light.

\paragraph{The first layer is the only one that multiplies.} Every other layer
receives binary spikes and only accumulates, but the first convolution receives
analog pixel values and performs full multiply-accumulates. Because its input is
identical at every step, so is its output current, which can be computed once
rather than $T$ times. Section~\ref{sec:energy-method} shows how this choice
 decides whether the spiking network is cheaper than its conventional
equivalent.

\paragraph{The input is deterministic.} Without sampling, two forward passes on
the same image produce identical spike trains, which is what allows two models to
be compared spike for spike on the same input, as the calibration of
Section~\ref{sec:calibration} requires.

\section{Quantization}
\label{sec:quantization-background}

Quantization represents the parameters and intermediate values of a network with
fewer bits than the 32 of floating point. Storing a value in $b$ bits divides its
memory by $32/b$, and integer operations on narrow operands cost far less energy
than floating-point ones (Section~\ref{sec:energy-background}), which is why
quantization is the standard practice adopted between a trained model and its deployment on
constrained hardware~\cite{jacob2018quantization, nagel2021whitepaper}.

\paragraph{The uniform quantizer.} A uniform quantizer maps a real value $x$ to
one of at most $2^b$ equally spaced levels: it divides by a step $\Delta$, rounds to the
nearest integer, clips to the representable range, and multiplies back,
\begin{equation}
\hat{x} = \Delta \cdot \mathrm{clip}\!\left(\mathrm{round}\!\left(\frac{x}{\Delta}\right),
\, q_{\min},\, q_{\max}\right)
\label{eq:quantizer}
\end{equation}

where $q_{\min}$ and $q_{\max}$ are the smallest and largest integer levels
representable with $b$ bits: $-(2^{b-1}-1)$ and $2^{b-1}-1$ for a symmetric
signed grid-- which keeps zero exactly representable at the cost of one unused
code-- $2^b - 1$ levels in all; $0$ and $2^b - 1$ for an unsigned one, which uses
all $2^b$. The largest representable value is therefore $\Delta\, q_{\max}$.

For a symmetric grid over $[-w_{\max}, +w_{\max}]$ the largest level must reach
$w_{\max}$, so the step is
\begin{equation}
\Delta = \frac{w_{\max}}{2^{b-1} - 1} = \frac{2 w_{\max}}{2^b - 2},
\label{eq:step}
\end{equation}
the step the weight quantizer of Chapter~\ref{ch:quantization} computes for each
output channel. Section~\ref{sec:perturbation} expresses a perturbation in a unit
derived from it: the 8-bit step of a single per-tensor grid set by the largest
weight of the whole network, computed there as $2 w_{\max}/(2^8 - 1)$, which is
$0.4\,\%$ smaller than Equation~\ref{eq:step} at 8 bits.
Since the quantizer
works per channel, and most channels have a smaller maximum than the network as
a whole, that unit is coarser than the steps the quantizer actually applies, and
the floor it defines is conservative.

There are two main errors that follow from Equation~\ref{eq:quantizer}, with the choice of range
trades one against the other: values inside the range are moved by at most
$\Delta/2$ (\emph{rounding} error), values outside it are cut to the nearest
bound (\emph{clipping} error). For values spread over many steps the rounding
error is close to uniform on $[-\Delta/2, \Delta/2]$, with a root-mean-square of
$\Delta/\sqrt{12}$; a Gaussian perturbation whose standard deviation is one step,
the unit of Section~\ref{sec:perturbation}, is therefore about $3.5$ times
larger than the rounding it stands for. A wide range avoids clipping but makes
$\Delta$ coarse; a narrow one refines $\Delta$ at the cost of saturating the
tails.

When the quantity is non-negative, or its negative values carry no information
for the task, an \emph{unsigned} grid over $[0, x_{\max}]$, with step
$x_{\max}/(2^b - 1)$, spends no levels on values the network does not use. The
membrane potential of this network is of the second kind. Negative potentials
are frequent --- in the first layer the median is $-0.08$ and the first
percentile $-6.6$ --- yet clipping them to zero moves the validation accuracy by
one image in $1\,435$ (Section~\ref{sec:membrane-prep}).

More specifically, the step can be shared by a whole tensor (\emph{per-tensor}) or computed
separately for each output channel (\emph{per-channel}). A per-tensor step is
set by the largest value in the tensor, so channels whose weights are
systematically smaller use only a fraction of the available levels; the
per-channel scheme avoids this at the price of one scale factor per
channel~\cite{krishnamoorthi2018whitepaper, nagel2021whitepaper}.

Given a quantizer, two families of method decide \emph{when} it is applied:
after training, or during it.

\paragraph{Post-training quantization (PTQ).} PTQ converts an already trained
full-precision model: the weights are rounded to their grid,
and the ranges of the activations, which depend on the data, are estimated on a
small unlabelled calibration set~\cite{nagel2021whitepaper}. It is a relatively fast method and
needs neither labelled data nor a training pipeline, which makes it the
``push-button'' option and the one most used for rapid deployment. Its limit is
precision: at 8 bits the rounding error is small compared with the spread of the
weights, and per-channel PTQ keeps classification accuracy within about 2\,\% of
floating point on standard convolutional networks. At 4 bits the step is sixteen
times coarser, and accuracy typically degrades unless more elaborate correction
methods are used~\cite{krishnamoorthi2018whitepaper, nagel2021whitepaper}. Rabbi et
al.~\cite{rabbi2026} use static weight-only PTQ at INT8 and INT4, the setting in which
they observe explanations degrading before accuracy.

\paragraph{Quantization-aware training (QAT).} QAT includes the quantizer in
training, so that the network can adapt its weights to the grid. Specifically, in the forward pass every quantized tensor is
passed through Equation~\ref{eq:quantizer} --- \emph{simulated} or
\emph{fake} quantization, since the values are rounded but still stored and
processed in floating point --- so that the loss is computed on the network as
it will be deployed~\cite{jacob2018quantization}. The difficulty lies in the backward
pass: rounding is a step function, whose derivative is zero almost everywhere,
and no gradient would reach the weights. The \emph{straight-through estimator}
solves it by treating rounding as the identity in the backward
pass~\cite{bengio2013ste}: the forward pass sees the quantized value, the
gradient flows as if it had not been quantized. This is the same device, in a
different guise, as the surrogate gradient of Section~\ref{sec:surrogate}, which
also keeps a step function in the forward pass and replaces its derivative in
the backward one. QAT needs labelled data and additional training, usually as a
fine-tuning of the full-precision model, and in exchange reaches lower
bit-widths: it narrows the 8-bit gap to about 1\,\% and makes 4-bit weights
usable, with losses between 2 and 10\,\% depending on the size of the
network~\cite{krishnamoorthi2018whitepaper, nagel2021whitepaper}.

The two families therefore answer different needs: PTQ when the model must be
compressed cheaply and 8 bits suffice, QAT when the target precision is lower
than PTQ can reach. In a spiking network both apply not only to the weights but
also to the membrane potential, as the next section describes.

\subsection{Quantization for spiking networks}

Spiking architectures add constraints of their own, as could already be inferred.
The membrane decay $\beta$
multiplies a stored state at every time step, so on integer hardware it is
cheapest when realizable as a bit shift: $\beta = 2^{-k}$ is a right shift, and
$\beta = 1 - 2^{-k}$ is a shift plus a subtraction. MINT~\cite{yin2024mint}
quantizes both targets uniformly and shares a single
scale factor between weights and membrane potentials. Since the weighted input
and the stored state are then expressed in the same integer units, they can be
added directly, and the multipliers that conventional uniform quantization needs
to rescale one into the other disappear from the inference loop. At 2 bits,
MINT reports $90.6\,\%$ accuracy on CIFAR-10 with a VGG-16 and a memory
reduction of $93.8\,\%$ with respect to full precision.
HardF-SNN~\cite{liu2026hardfsnn} shows the cost of that sharing: weights have a
narrow, centred distribution, whereas membrane potentials have a much broader
one with a long tail, so a scale fitted to the weights compresses the dynamic
range of the membrane and collapses many distinct potentials onto the same
level. It therefore scales the membrane in proportion to the weights rather
than identically, replaces batch normalization with an integer, shift-based
version, and verifies the result on an FPGA accelerator, noting that most
quantized spiking networks are still evaluated with simulated quantization, in
which part of the computation remains in floating point.

Smith et al.~\cite{smith2026} argue independently that accuracy is an
insufficient criterion for quantized spiking networks. They in fact propose an Earth
Mover's Distance between per-neuron firing-rate distributions as a diagnostic of
behavioural divergence. Separating weight from membrane quantization as this
thesis does, and reporting that membrane-potential quantization requires at least
four bits to preserve both accuracy and dynamics while learned weight
quantization remains viable at two. They ground the asymmetry mechanistically: a
weight quantization error is fixed for an entire forward pass, whereas a membrane
error propagates through the recurrent update at every time step.

This thesis reaches the opposite result at 2 bits: the membrane arm keeps
accuracy at parity (Section~\ref{sec:grid-results}). A plausible reason is where
the quantizer acts. Here it rounds the state carried into the next step, after
the spike decision (Section~\ref{sec:membrane-quantizer}). At every step the
threshold is compared with $\beta$ times the quantized previous state plus a
full-precision input current, so the comparison itself is never made on a 2-bit
value.

\subsection{Energy and memory accounting}
\label{sec:energy-background}

Since no model is deployed on hardware in this work, energy is estimated
analytically, as is standard in the spiking literature: the operations an
inference requires are counted and multiplied by a per-operation cost at a
reference process node, conventionally the 45\,nm figures of
Horowitz~\cite{horowitz2014}. The result is a
like-for-like comparison between configurations of the same network.

\paragraph{Counting operations.} In a conventional layer every connection
performs one multiply-accumulate (MAC) per inference. In a spiking layer the
input is binary, so a connection only \emph{adds} its weight, and only when the
presynaptic neuron fires. The number of these synaptic operations (SOP) is
therefore the dense MAC count scaled by the firing rate $f$ of the input layer
and by the number of time steps,
\begin{equation}
\mathrm{SOP} = f \cdot T \cdot \mathrm{MAC}.
\end{equation}
In 32-bit floating point a MAC costs $4.6$\,pJ (a $3.7$\,pJ multiplication plus a
$0.9$\,pJ addition) and an accumulate $0.9$\,pJ. That factor of five is the
physical reason a spiking network can be cheaper and lighter, and it also sets the limit:
a spiking layer saves energy only while $f \cdot T < 4.6 / 0.9 \approx 5$, so
sparsity and a short time window are both necessary. The first layer is an
exception, since under direct encoding it receives analog values and performs
full MACs (Section~\ref{sec:encoding}).

\paragraph{Quantized operations.} The same table supplies the integer figures
used for quantized operating points: an 8-bit multiplication costs $0.2$\,pJ and
an 8-bit addition $0.03$\,pJ, thirty times less than its floating-point
counterpart, so that an 8-bit MAC costs $0.23$\,pJ under the same convention
that gives $4.6$\,pJ in floating point. Narrower integer operations are not tabulated, so any figure below
eight bits is an extrapolation, whose rule is stated in
Chapter~\ref{ch:quantization}.

\paragraph{What the count leaves out.} The estimate covers arithmetic only. It
ignores the update of each neuron (leak, threshold comparison, reset), control
logic, and above all data movement, which on real hardware is often the larger
cost: in the same 45\,nm table, reading 64 bits from on-chip SRAM costs between
$10$ and $100$\,pJ depending on its size, and from off-chip DRAM between $1.3$
and $2.6$\,nJ~\cite{horowitz2014}, one to three orders of magnitude more than an
arithmetic operation.

\paragraph{Why memory is accounted separately.} For this reason the memory
footprint is reported alongside the energy estimate:
the energy of an access depends on where the data lives, which is a property of
the hardware, whereas the number of bits to be stored is a property of the
model. The footprint is computed as the bits needed for the weights plus the
bits needed for the membrane state of every neuron. The second term is specific
to spiking networks: a conventional network can discard the activations of a
layer once the next one has been computed, but a spiking network must keep the
potential of every neuron from one time step to the next for the whole
inference. That state scales with the number of neurons rather than the number
of parameters, and it can be quantized independently of the weights
(Section~\ref{sec:quantization-background}).

\section{Explainability for spiking networks}
\label{sec:sam-background}

Attribution methods developed for conventional networks transfer badly to
spiking ones, for two different causes. Gradient-based methods such as
Grad-CAM~\cite{selvaraju2017gradcam} weight feature maps by the gradient of the
class score, and in a spiking network that gradient exists only as a surrogate
(Section~\ref{sec:surrogate}). The approximation is applied again at every layer
the gradient crosses, so the error accumulates with depth: Kim and
Panda~\cite{kim2021sam} observe that Grad-CAM transplanted to spiking networks
produces smoothed, poorly discriminative maps, worst at the shallow layers.
Perturbation methods such as LIME~\cite{ribeiro2016lime} avoid gradients
altogether and are architecture-agnostic, but they need hundreds to thousands of
forward passes per explanation. Both families, moreover, return a single map per input,
and so discard the temporal axis along which a spiking network computes.

\paragraph{The Spike Activation Map.} The Spike Activation Map (SAM) of Kim and
Panda~\cite{kim2021sam} was designed for spiking networks from the start. It
rests on a biological observation: spikes that follow each other at short
inter-spike intervals are more likely to drive the membrane of the downstream
neuron above threshold, and therefore carry more information than isolated
ones. SAM numerically quantifies this into a score computed in three steps. A \emph{temporal
spike contribution score} weights a past spike at $t'$ by its recency with
respect to the current step $t$,
\begin{equation}
T(t, t') = \exp\left(-\gamma\,|t - t'|\right),
\end{equation}
where $\gamma$ sets how fast the contribution of a spike fades. A
\emph{neuronal contribution score} sums these weights over the spikes the neuron
emitted before $t$,
\begin{equation}
N^{k}_{ij,t} = \sum_{t' \in P^{k}_{ij},\; t' < t} T(t, t'),
\end{equation}
where $P^{k}_{ij}$ is the set of firing times of the neuron in channel $k$ at
position $(i,j)$; it is high for a neuron that has fired often and recently, and
low for one that fired rarely or long ago. Finally, the map at time $t$ keeps
the score only of the neurons that are firing at $t$ and sums over the channels
of the layer,
\begin{equation}
M_{ij,t} = \sum_{k} N^{k}_{ij,t}\, S^{k}_{ij,t},
\end{equation}
yielding one spatial heat map per time step and per layer.

\paragraph{Properties.} Several properties of SAM matter for this thesis. It uses
solely the spike trains of the forward pass, so the explanation describes the function the
network actually computes, and its cost is marginal --- a store of recent spike
times and a lookup table for the kernel, which Kim and Panda note could be
implemented in an inference accelerator. It needs no class label: SAM shows
where the network concentrates its activity for a given input, not the evidence
for a particular class, so its faithfulness must be tested against the class
the model predicts. And it has a temporal axis: an explanation is a sequence of
$T$ maps, which is what makes it possible to ask whether quantization changes
\emph{when} a network decides as well as \emph{where} it looks. Kim and Panda
validate the method on Tiny-ImageNet, finding SAM more discriminative than
Grad-CAM adapted to spiking networks.

Feature-attribution methods for spiking networks have also been surveyed more
broadly~\cite{nguyen2023attribution}, but SAM remained the method that reads the
spike trains directly, and it is the explainer used throughout this work. Its
one free parameter, $\gamma$, is fixed in Section~\ref{sec:gamma}.

\subsection{Evaluating an explanation}

A stable explanation can still be wrong, so stability requires a separate test
of faithfulness. The \emph{deletion metric} of Petsiuk et
al.~\cite{petsiuk2018rise} removes pixels in decreasing order of attributed
importance, and records how quickly the model's confidence falls, summarized as
the area under that curve. The faster the fall, the smaller the area. The
underlying hypothesis is that a faithful map degrades the prediction faster than
a map that carries no information. Petsiuk et al. read the area on its own; this
thesis compares it with the removal of pixels in random order on the same image
(Section~\ref{sec:faithfulness}), so that a map counts as faithful only by the
margin it gains over chance.

The metrics used to quantitatively compare two explanations with each other are borrowed from
the saliency-model literature, where they were specifically developed to compare a predicted
map against human eye fixations~\cite{bylinskii2019metrics}: Pearson
correlation, structural similarity~\cite{wang2004ssim}, and the
intersection-over-union of the highest-scoring regions. Bylinskii et al. observe
that rankings produced by different metrics frequently disagree, which is the
main reason for reporting more than one. This thesis reports the correlation and
the overlap. Structural similarity, which Rabbi et al.~\cite{rabbi2026} use
alongside them, is left out as redundant with the two.

A separate line of work asks whether attribution methods are testing anything at
all. Adebayo et al.~\cite{adebayo2018sanity} randomize model weights
progressively and verify that explanations change in response, as a map that
survives the destruction of the model it sets to explain is not explaining
it. The perturbation experiment of Section~\ref{sec:perturbation} is a graded
form of that test, run in the opposite direction.

\section{Gravitational-wave glitch classification}
\label{sec:gravityspy}

The sensitivity of the LIGO interferometers is primarily limited by non-Gaussian transient
noise of instrumental and environmental origin. As introduced in
Chapter~\ref{ch:introduction}, these glitches are short, can mimic astrophysical
signals, and cna be grouped into morphological families that correlate with
distinct physical causes.

The Gravity Spy project~\cite{zevin2017, bahaadini2018} combines machine
classification with citizen science to label them. The release used
here~\cite{glanzer2021} covers observing runs O1 through O3b and provides, for
each Omicron trigger, the trigger metadata, the per-class scores of a
convolutional classifier, the label that classifier assigns, and the addresses
of the Omega-scan spectrograms, which are downloaded separately from the
Zooniverse archive. The labels used in this work are therefore machine labels,
not volunteer consensus, and no threshold on the classifier's confidence was
applied when the data were selected.

The scores released with the labels do not always support them. Of the
$9\,568$ triggers selected for the four classes, the score of the assigned class
is below $0.5$ for $24\,\%$ and below $0.01$ for $615$, and for $4.4\,\%$ the
label is not the class with the highest of the four scores. The disagreement is
concentrated in two classes: the score falls below $0.5$ for $53\,\%$ of the
Violin\_Mode triggers and $37\,\%$ of the Extremely\_Loud ones, against $4\,\%$
of Blip and $3\,\%$ of Scattered\_Light. An accuracy measured against these
labels is an agreement with the Gravity Spy classifier, not with a verified
truth.

The classification task is therefore image classification over time--frequency
representations. Its interest for this thesis is not its difficulty --- at
$98\,\%$ agreement with the reference classifier it is not difficult --- but
that the explanations have a consumer. Glitch classification is directly
connected with detector characterization, in which the
morphology of a glitch family is traced back to an instrumental cause and the
instrument is modified accordingly. An explanation that points at the wrong part of the
spectrogram sends a commissioner to the wrong causes.

\section{The gap this thesis occupies}
\label{sec:gap}

There are two close works useful to frame this thesis. Rabbi et
al.~\cite{rabbi2026} measure explanation fidelity under quantization, but for
conventional networks, implementing post-training weight-only quantization, and relying purely
on spatial explanations. Smith et al.~\cite{smith2026} address spiking networks and
separate the two quantization targets, but measure divergence in \emph{population
firing statistics} aggregated over a dataset, and do not take into account explanations.

The unoccupied ground is therefore a per-sample, spatially and temporally resolved
explanation fidelity for a spiking classifier with the two quantization targets
separated.